\section{Dataset}
\label{sec:dataset}

We evaluate the monitoring formulation on a controlled synthetic dataset. The synthetic design is deliberate: the research question requires the same underlying operational scenario to be rendered under different transcript-structure conditions while keeping task content fixed. A source-grounded synthetic dataset makes this possible, allowing differences between transcript conditions to be attributed to representation rather than scenario content.

\subsection{Source-Grounded Construction}

The dataset was constructed with a prompt-based generation and validation pipeline grounded in three source types: Swiss radio-procedure material\footnote{Bundesamt für Bevölkerungsschutz, Funkmaterial Sprechregeln, February 2004}, a real firefighter communication transcript used for surface realism, and canonical procedural regulations\footnote{Feuerwehr Koordination Schweiz FKS, Reglement Einsatzführung, November 2022}. These materials were distilled into source notes used for both scenario generation and validation.

Each scenario was generated as a fixed-schema JSON object containing an ordered message sequence, a closed-world list of predefined tasks, and gold task states. In addition to the final task state, each gold entry stores the message id at which assignment first becomes valid and, when applicable, the message id at which completion first becomes valid. Task-command linkage is annotation-grounded rather than based on lexical overlap between task labels and command wording. This enables deterministic derivation of gold states for every incremental transcript prefix.

Validation was performed in two independent rubric-guided passes: a structural pass for radio style and protocol coherence, and a content pass for operational plausibility, role-task alignment, task sequencing, and traceability of gold states to explicit message evidence. We implemented both dataset generation and validation with \texttt{gpt-5.4} as a bounded LLM judgment process under source-grounded constraints \citep{Zheng2023,Rotteger2024SafetyPrompts}. This source-grounded generation-and-judge setup was used to make the synthetic data more representative of the intended application setting while preserving control over transcript-structure variation. The generation--validation--revision loop ran for 12 rounds until the dataset passed both validation streams.

As an additional plausibility check, the extracted structure and resulting scenarios were reviewed within the team. They were judged to align with observed communication structures and contents from the intended application setting. While informal and not a formal annotation study, this review provides additional support for the dataset's plausibility as a set of realistic scenarios.

\subsection{Dataset Contents}

The final dataset contains five German-language firefighter scenarios with 102 ordered radio messages and 15 predefined tasks in total. Every scenario contains exactly three monitored tasks, yielding 15 task-state traces overall. All 15 tasks are explicitly assigned in the transcript evidence, 12 are completed within the scenario, and 3 remain assigned but incomplete at scenario end.

The five scenarios cover a kitchen fire in an apartment building, an underground-garage fire, a school basement fire, a workshop fire involving gas cylinders, and an attic fire in a row house. Across these scenarios, the monitored task families include water supply, search, ventilation, suppression, access control, gas-cylinder cooling, CO measurement, and final control. The assigned units include command-relevant firefighter roles such as \texttt{Angriffstrupp}, \texttt{Wassertrupp}, \texttt{Messgruppe}, \texttt{Kontrolltrupp}, and \texttt{Verkehrstrupp}.

\subsection{Representation and Limitations}

Each scenario stores ordered messages, predefined tasks, and one final gold state per task, using the same output fields as the monitoring model: assignment status, assigned unit, completion status, and optional completion outcome. For incremental evaluation, per-prefix gold states are derived deterministically from the stored assignment and completion transition ids.

These gold task states serve as the reference labels for evaluation, and the stored transition ids define the gold assignment and completion points used in incremental detection metrics.
For evaluation, these fields also induce the three-way task state \texttt{NOT\_ASSIGNED}/\texttt{ASSIGNED}/\texttt{COMPLETED}.

Before evaluation, all scenario files pass deterministic local validation checks for schema conformance, sequential message ids, exact alignment between predefined tasks and gold task states, and consistency constraints on assignment/completion labels.

The limitations of the study follow directly from this dataset design and should be read upfront. The dataset is synthetic, small, and intentionally controlled. The scenarios are cleaner and more linear than real firefighter radio traffic, with limited protocol noise, dispatch and handover tails, and some bundled tasks that compress multi-step work into a single monitored item. The evaluation also begins from text transcripts rather than raw audio, so it does not include upstream speech-recognition or radio-channel degradation effects. Finally, the task formulation is deliberately narrow: it tests closed-world monitoring of predefined tasks, not open-ended task discovery or broad operational performance. We therefore use the dataset as controlled evidence for transcript-structure sensitivity in dashboard-oriented task monitoring, not as a realism corpus or a direct deployment claim.
